% Source: https://artofproblemsolving.com/community/c3h1521792_help_please
% Stable identity: aops-topic-1521792-q01
Magnifying by a factor of $6$ is a dilation: every length is multiplied by $6$, so every triangle is mapped to a similar triangle (its sides are in the same ratios). Similar triangles have equal angles, so a dilation does not change angle measures. The sides of the angle look $6$ times as long, but the opening between them is the same (in the sketch the angle is drawn wider than $6^\circ$ for clarity):

\begin{asy}
usepackage("amsmath");
size(7cm);
defaultpen(fontsize(10pt));
pair O=(0,0);
real t=28;
draw(O--(0.5,0),linewidth(1));draw(O--0.5*dir(t),linewidth(1));
pair V=(1.3,0);
draw(V--V+(3,0),blue+linewidth(1));draw(V--V+3*dir(t),blue+linewidth(1));
draw(arc(O,0.3,0,t),red);label("$6^\circ$",0.42*dir(t/2),E,red+fontsize(8pt));
draw(arc(V,1.2,0,t),red);label("$6^\circ$",V+1.25*dir(t/2),E,red);
label("original",(0.25,-0.1),S,fontsize(8pt));
label("seen through the glass ($6\times$ longer sides)",V+(1.5,-0.1),S,fontsize(8pt)+blue);
label("\mbox{}",(min(currentpicture,user=true)+max(currentpicture,user=true))/2,black); // keeps the text after the figure black when LaTeX typesets the labels inline
\end{asy}

The angle still measures $\boxed{6^\circ}$.
